% !TeX program = lualatex
% ================================================================
% Polish Parallel Text Version of IterativeProcessesWorksheet
%
% Generated by the server using the following settings:
%
%   language            Polish (pol)
%   text direction      left to right
%   font                Noto Serif
%   fallback font       Noto Serif
%   built from          latex/worksheets/percentages/iterativeProcessesWorksheet.tex
%   vocabulary key      latex/worksheets/percentages/iterativeProcessesWorksheet/L2/pol/iterativeProcessesWorksheet_polishKey.tex
%   tier 3 vocabulary   green   RGB 0 128 0
%   English text        black   RGB 0 0 0
%   Polish text         purple  RGB 112 48 160
%   layout macros       version 3
%
% Please don't edit any information in this comment block - the server
% compares it against the current settings and rebuilds this file when
% they differ, so changes here would be overwritten. However, feel free
% to make updates to the LaTeX below if you want to edit it.
% ================================================================

\documentclass[12pt,a4paper]{article}
% ================================================================
% Copied from the English sheet this was translated from, so that
% anything it sets up for itself still works here
% ================================================================
% ── Packages ────────────────────────────────────────────────────────────────
\usepackage[a4paper, top=2cm, bottom=2cm, left=2.2cm, right=2.2cm]{geometry}
\usepackage{amsmath, amssymb}
\usepackage{booktabs}
\usepackage{array}
\usepackage{enumitem}
\usepackage{tcolorbox}
\tcbuselibrary{skins, breakable}
\usepackage{multicol}
\usepackage{tabularx}
\usepackage{colortbl}
\usepackage{parskip}
\usepackage[T1]{fontenc}

% ── Colours ─────────────────────────────────────────────────────────────────
\definecolor{headgrey}{gray}{0.15}
\definecolor{ruleblue}{RGB}{30,80,160}
\definecolor{scaffoldgold}{RGB}{255,248,220}
\definecolor{scaffoldborder}{RGB}{200,160,40}
\definecolor{sectionbg}{RGB}{240,244,252}
\definecolor{sectionrule}{RGB}{30,80,160}
\definecolor{challengebg}{RGB}{245,240,255}
\definecolor{challengerule}{RGB}{100,60,180}
\definecolor{answerbg}{gray}{0.97}
\definecolor{lightgrey}{gray}{0.88}

% ── Section header box ───────────────────────────────────────────────────────
\newcommand{\sectionbox}[2]{%
  \begin{tcolorbox}[
    colback=sectionbg, colframe=sectionrule,
    leftrule=4pt, rightrule=0pt, toprule=0pt, bottomrule=0pt,
    sharp corners, left=6pt, top=4pt, bottom=4pt
  ]
  {\large\bfseries\color{sectionrule} #1}\quad{\normalsize #2}
  \end{tcolorbox}%
}

% ── Scaffold box ─────────────────────────────────────────────────────────────
\newcommand{\scaffold}[1]{%
  \begin{tcolorbox}[
    colback=scaffoldgold, colframe=scaffoldborder,
    leftrule=3pt, rightrule=0pt, toprule=0pt, bottomrule=0pt,
    sharp corners, left=6pt, top=3pt, bottom=3pt, before skip=4pt
  ]
  {\small\itshape\color{brown} \textbf{Getting Started:} #1}
  \end{tcolorbox}%
}

% ── Answer lines ─────────────────────────────────────────────────────────────
\newcommand{\answerlines}[1]{%
  \vspace{2pt}%
  \foreach \i in {1,...,#1}{%
    \vspace{0.55cm}\noindent\textcolor{lightgrey}{\rule{\linewidth}{0.4pt}}\par
  }%
  \vspace{4pt}%
}

% ── Answer box (blank area) ───────────────────────────────────────────────────
\newcommand{\answerbox}[1]{%
  \begin{tcolorbox}[
    colback=answerbg, colframe=lightgrey,
    sharp corners, left=4pt, top=4pt, bottom=4pt,
    height=#1, valign=top
  ]
  \end{tcolorbox}%
}

% ── Misc ─────────────────────────────────────────────────────────────────────
\newcolumntype{G}{>{\columncolor{sectionbg}}c}
\setlist[enumerate,1]{label=(\alph*), leftmargin=*, itemsep=0.4em}
\setlist[enumerate,2]{label=\roman*., leftmargin=2em}
\pagestyle{plain}

% ════════════════════════════════════════════════════════════════════════════

% Characters this language's font has not got are borrowed from another,
% rather than being silently dropped - see docs/translated-sheet-latex.md
\directlua{%
  if luaotfload and luaotfload.add_fallback then
    luaotfload.add_fallback("eallatinfallback",
      { "Noto Serif:mode=harf;" })
  else
    tex.error("this luaotfload is too old to register a fallback font")
  end
}
\usepackage[english]{babel}
\babelprovide[import]{polish}
\babelfont[polish]{rm}[Renderer=Harfbuzz, BoldFont={Noto Serif}, ItalicFont={Noto Serif}, BoldItalicFont={Noto Serif}, RawFeature={fallback=eallatinfallback}]{Noto Serif}
\babelfont[polish]{sf}[Renderer=Harfbuzz, BoldFont={Noto Serif}, ItalicFont={Noto Serif}, BoldItalicFont={Noto Serif}, RawFeature={fallback=eallatinfallback}]{Noto Serif}
\newcommand{\ealtext}[1]{\foreignlanguage{polish}{#1}}
\newcommand{\ealtextblock}[1]{{\raggedright\foreignlanguage{polish}{#1}\par}}
% ================================================================
% EAL layout helpers
% ================================================================
\usepackage{xcolor}

\definecolor{ealkeycolour}{RGB}{0,128,0}
\definecolor{ealenglish}{RGB}{0,0,0}
\definecolor{ealtwocolour}{RGB}{112,48,160}

\newsavebox{\ealboxen}
\newsavebox{\ealboxtr}
\newlength{\ealglosswd}
\newlength{\ealglossraise}
\setlength{\ealglossraise}{1.15em}

% a tier 3 word, in English and in translation alike
\newcommand{\ealkey}[1]{\textcolor{ealkeycolour}{#1}}
\newcommand{\ealkeytr}[1]{\textcolor{ealkeycolour}{\ealtext{#1}}}

% One translated word above one English word. Built from kernel
% primitives, never a tabular, which would break wherever the sheet
% loads array, booktabs or tabularx. The box is as wide as the wider of
% the two, so a long translation pushes its neighbours aside instead of
% overlapping them, and it claims its full height so a table row or a
% TikZ node makes room for it.
\newcommand{\ealgloss}[2]{%
  \sbox{\ealboxen}{\ealkey{#1}}%
  \sbox{\ealboxtr}{\small\textcolor{ealtwocolour}{\ealtext{#2}}}%
  \setlength{\ealglosswd}{\wd\ealboxen}%
  \ifdim\wd\ealboxtr>\ealglosswd\setlength{\ealglosswd}{\wd\ealboxtr}\fi
  \leavevmode
  \makebox[\ealglosswd][c]{%
    \makebox[0pt][c]{\usebox{\ealboxen}}%
    \makebox[0pt][c]{\raisebox{\ealglossraise}{\usebox{\ealboxtr}}}%
  }%
}

% opens up line spacing around a block of glosses, so the raised
% translations sit clear of the line above
\newenvironment{ealglossed}{\par\linespread{2.1}\selectfont\ignorespaces}{\par}

% a whole translated sentence above its English counterpart. block level,
% so it does not belong inside a tabular cell
\newcommand{\ealpara}[2]{%
  \par\addvspace{0.5\baselineskip}%
  {\color{ealtwocolour}\ealtextblock{#1}}%
  \penalty200
  {\color{ealenglish}#2\par}%
  \addvspace{0.5\baselineskip}%
}

\begin{document}
% ════════════════════════════════════════════════════════════════════════════

% ── Title block ─────────────────────────────────────────────────────────────
\begin{center}
  {\LARGE\bfseries\color{headgrey} \ealpara{Procesy iteracyjne}{\ealkey{Iterative Processes}}}
  \vspace{4pt}
  {\large \ealpara{Procesy iteracyjne \textbullet\ Tabele wartości \textbullet\ \ealkeytr{Wzory na $n$-ty wyraz}}{Iterative Processes \textbullet\ Tables of Values \textbullet\ \ealkey{$n$th Terms}}}
  \vspace{8pt}
  \textcolor{lightgrey}{\rule{\linewidth}{1pt}}
\end{center}

\vspace{-2pt}
\noindent
\begin{tabularx}{\linewidth}{X X X}
  \textbf{Name:}\enspace\dotfill & \textbf{Class:}\enspace\dotfill & \textbf{Date:}\enspace\dotfill
\end{tabularx}

\vspace{10pt}

% ── Reference panel ─────────────────────────────────────────────────────────
\begin{tcolorbox}[
  colback=white, colframe=headgrey,
  sharp corners, title={\bfseries Key Ideas},
  fonttitle=\small\bfseries, coltitle=white, colbacktitle=headgrey,
  top=6pt, bottom=6pt, left=8pt, right=8pt
]
\small
\begin{multicols}{2}
\ealpara{\textbf{\ealkeytr{Proces iteracyjny}} łączy każdy \ealkeytr{wyraz} z poprzednim.}{\textbf{An \ealkey{iterative process}} links each \ealkey{term} to the one before it.}

\ealpara{Przykład: $\quad U_{n+1} = U_n + 3, \quad U_1 = 5$}{Example: $\quad U_{n+1} = U_n + 3, \quad U_1 = 5$}

\ealpara{To oznacza: \emph{dodaj 3 do poprzedniego \ealkeytr{wyrazu}, zaczynając od 5.}}{This means: \emph{add 3 to the previous \ealkey{term}, starting at 5.}}

\ealpara{Wyrazy: $5,\ 8,\ 11,\ 14, \ldots$}{Terms: $5,\ 8,\ 11,\ 14, \ldots$}

\columnbreak

\ealpara{\ealkeytr{Wzór na $n$-ty wyraz} daje dowolny \ealkeytr{wyraz} bezpośrednio.}{\textbf{An \ealkey{$n$th term} formula} gives any \ealkey{term} directly.}

\ealpara{Przykład: $\quad U_n = 3n + 2$}{Example: $\quad U_n = 3n + 2$}

\ealpara{To oznacza: \emph{\ealkeytr{podstaw} dowolne $n$, aby natychmiast otrzymać ten \ealkeytr{wyraz}.}}{This means: \emph{\ealkey{substitute} any $n$ to get that \ealkey{term} immediately.}}
\end{multicols}
\end{tcolorbox}

\vspace{14pt}

% ════════════════════════════════════════════════════════════════════════════
\sectionbox{\ealpara{Sekcja A}{Section A}}{\ealpara{\ealkeytr{Proces iteracyjny} $\longrightarrow$ Tabela wartości}{\ealkey{Iterative Process} $\longrightarrow$ Table of Values}}
% ════════════════════════════════════════════════════════════════════════════

\bigskip

% ── A1 ──────────────────────────────────────────────────────────────────────
\ealpara{\textbf{A1.}\quad Rozważ \ealkeytr{proces iteracyjny}}{\textbf{A1.}\quad Consider the \ealkey{iterative process}}
\[
  U_{n+1} = U_n + 7, \qquad U_1 = 4
\]

\scaffold{\ealpara{Aby znaleźć $U_2$, \ealkeytr{podstaw} $n=1$: \enspace $U_2 = U_1 + 7 = 4 + 7 = 11$. Teraz kontynuuj w ten sam sposób dla każdego kolejnego \ealkeytr{wyrazu}.}{To find $U_2$, \ealkey{substitute} $n=1$: \enspace $U_2 = U_1 + 7 = 4 + 7 = 11$. Now continue in the same way for each subsequent \ealkey{term}.}}

\begin{enumerate}
  \item \ealpara{Uzupełnij poniższą tabelę wartości.}{Complete the table of values below.}

  \vspace{6pt}
  \begin{center}
  \renewcommand{\arraystretch}{1.8}
  \begin{tabular}{|>{\bfseries\centering\arraybackslash}m{1.2cm}
                  |>{\centering\arraybackslash}m{1.6cm}
                  |>{\centering\arraybackslash}m{1.6cm}
                  |>{\centering\arraybackslash}m{1.6cm}
                  |>{\centering\arraybackslash}m{1.6cm}
                  |>{\centering\arraybackslash}m{1.6cm}
                  |>{\centering\arraybackslash}m{1.6cm}|}
  \hline
  \rowcolor{sectionbg}
  $n$     & 1 & 2 & 3 & 4 & 5 & 6 \\
  \hline
  $U_n$   & 4 &   &   &   &   &   \\
  \hline
  \end{tabular}
  \end{center}

  \vspace{8pt}
  \item \ealpara{Opisz słowami, co robi ten \ealkeytr{ciąg}.}{Describe in words what this \ealkey{sequence} is doing.}

  \answerlines{2}

\end{enumerate}

\bigskip\bigskip

\newpage{}

% ── A2 ──────────────────────────────────────────────────────────────────────
\ealpara{\textbf{A2.}\quad Rozważ \ealkeytr{proces iteracyjny}}{\textbf{A2.}\quad Consider the \ealkey{iterative process}}
\[
  V_{k+1} = 3V_k, \qquad V_1 = 2
\]

\scaffold{\ealpara{Operacją tutaj jest \emph{mnożenie}, a nie dodawanie. Aby znaleźć $V_2$: \enspace $V_2 = 3 \times V_1 = 3 \times 2 = 6$. Zauważ, że \ealkeytr{zmienną} jest $V_k$, a nie $U_n$ --- proces działa w dokładnie ten sam sposób.}{The operation here is \emph{multiplication}, not addition. To find $V_2$: \enspace $V_2 = 3 \times V_1 = 3 \times 2 = 6$. Notice the \ealkey{variable} is $V_k$, not $U_n$ --- the process works in exactly the same way.}}

\begin{enumerate}
  \item \ealpara{Uzupełnij poniższą tabelę wartości.}{Complete the table of values below.}

  \vspace{6pt}
  \begin{center}
  \renewcommand{\arraystretch}{1.8}
  \begin{tabular}{|>{\bfseries\centering\arraybackslash}m{1.2cm}
                  |>{\centering\arraybackslash}m{1.6cm}
                  |>{\centering\arraybackslash}m{1.6cm}
                  |>{\centering\arraybackslash}m{1.6cm}
                  |>{\centering\arraybackslash}m{1.6cm}
                  |>{\centering\arraybackslash}m{1.6cm}|}
  \hline
  \rowcolor{sectionbg}
  $k$     & 1 & 2 & 3 & 4 & 5 \\
  \hline
  $V_k$   & 2 &   &   &   &   \\
  \hline
  \end{tabular}
  \end{center}

  \vspace{8pt}
  \item \ealpara{Opisz słowami, co robi ten \ealkeytr{ciąg}.}{Describe in word what the \ealkey{sequence} is doing.}

  \answerlines{2}

\end{enumerate}

\bigskip\bigskip

% ── A3 ──────────────────────────────────────────────────────────────────────
\ealpara{\textbf{A3.}\quad Rozważ \ealkeytr{proces iteracyjny}}{\textbf{A3.}\quad Consider the \ealkey{iterative process}}
\[
  F_{j+1} = F_j - 9, \qquad F_0 = 72
\]

\scaffold{\ealpara{Ten \ealkeytr{ciąg} \emph{zaczyna się od $j = 0$}, a nie $j = 1$. Pierwszym wpisem w twojej tabeli jest zatem $F_0 = 72$. Znajdź $F_1 = F_0 - 9 = 72 - 9 = 63$ i kontynuuj od tego miejsca.}{This \ealkey{sequence} \emph{starts at $j = 0$}, not $j = 1$. The first entry in your table is therefore $F_0 = 72$. Find $F_1 = F_0 - 9 = 72 - 9 = 63$, and continue from there.}}

\begin{enumerate}
  \item \ealpara{Uzupełnij poniższą tabelę wartości.}{Complete the table of values below.}

  \vspace{6pt}
  \begin{center}
  \renewcommand{\arraystretch}{1.8}
  \begin{tabular}{|>{\bfseries\centering\arraybackslash}m{1.2cm}
                  |>{\centering\arraybackslash}m{1.6cm}
                  |>{\centering\arraybackslash}m{1.6cm}
                  |>{\centering\arraybackslash}m{1.6cm}
                  |>{\centering\arraybackslash}m{1.6cm}
                  |>{\centering\arraybackslash}m{1.6cm}
                  |>{\centering\arraybackslash}m{1.6cm}
                  |>{\centering\arraybackslash}m{1.6cm}|}
  \hline
  \rowcolor{sectionbg}
  $j$     & 0  & 1 & 2 & 3 & 4 & 5 & 6 \\
  \hline
  $F_j$   & 72 &   &   &   &   &   &   \\
  \hline
  \end{tabular}
  \end{center}

  \vspace{8pt}
  \item \ealpara{Dla jakiej wartości $j$ $F_j$ po raz pierwszy staje się ujemne? Pokaż, skąd to wiesz.}{For what value of $j$ does $F_j$ first become negative? Show how you know.}

  \answerlines{3}


\end{enumerate}

\bigskip\bigskip

\newpage{}

% ── A4 ──────────────────────────────────────────────────────────────────────
\ealpara{\textbf{A4.}\quad Rozważ \ealkeytr{proces iteracyjny}}{\textbf{A4.}\quad Consider the \ealkey{iterative process}}
\[
  P_{n+1} = \frac{P_n}{4}, \qquad P_1 = 1024
\]

\scaffold{\ealpara{Dzielenie przez 4 to to samo, co mnożenie przez $\tfrac{1}{4}$. Początek: $P_1 = 1024$, więc $P_2 = \tfrac{1024}{4} = 256$. Kontynuuj, aż otrzymasz 5 \ealkeytr{wyrazów}.}{Dividing by 4 is the same as multiplying by $\tfrac{1}{4}$. Start: $P_1 = 1024$, so $P_2 = \tfrac{1024}{4} = 256$. Continue until you have 5 \ealkey{terms}.}}

\begin{enumerate}
  \item \ealpara{Uzupełnij poniższą tabelę wartości.}{Complete the table of values below.}

  \vspace{6pt}
  \begin{center}
  \renewcommand{\arraystretch}{1.8}
  \begin{tabular}{|>{\bfseries\centering\arraybackslash}m{1.2cm}
                  |>{\centering\arraybackslash}m{1.6cm}
                  |>{\centering\arraybackslash}m{1.6cm}
                  |>{\centering\arraybackslash}m{1.6cm}
                  |>{\centering\arraybackslash}m{1.6cm}
                  |>{\centering\arraybackslash}m{1.6cm}|}
  \hline
  \rowcolor{sectionbg}
  $n$     & 1    & 2 & 3 & 4 & 5 \\
  \hline
  $P_n$   & 1024 &   &   &   &   \\
  \hline
  \end{tabular}
  \end{center}

  \vspace{8pt}
  \item \ealpara{Zapisz każdą wartość $P_n$ jako \ealkeytr{potęgę} 4. Pierwsze dwie są już zrobione.}{Write each value of $P_n$ as a \ealkey{power} of 4. The first two are done for you.}

  \vspace{6pt}
  \begin{center}
  \renewcommand{\arraystretch}{1.8}
  \begin{tabular}{|>{\bfseries\centering\arraybackslash}m{1.2cm}
                  |>{\centering\arraybackslash}m{1.6cm}
                  |>{\centering\arraybackslash}m{1.6cm}
                  |>{\centering\arraybackslash}m{1.6cm}
                  |>{\centering\arraybackslash}m{1.6cm}
                  |>{\centering\arraybackslash}m{1.6cm}|}
  \hline
  \rowcolor{sectionbg}
  $n$     & 1     & 2     & 3 & 4 & 5 \\
  \hline
  $P_n$   & $4^5$ & $4^4$ &   &   &   \\
  \hline
  \end{tabular}
  \end{center}

  \vspace{8pt}
  \item \ealpara{Użyj schematu, aby zapisać \ealkeytr{wzór na $n$-ty wyraz} dla $P_n$ bezpośrednio.}{Use the pattern to write an \ealkey{$n$th term} formula for $P_n$ directly.}

 \vspace{1cm}

  \[
    P_n = \hspace{6cm}
  \]

  \vspace{1cm}

  \item \ealpara{Zapisz $P_n$ w postaci $P_n = \underline{\phantom{4096}} \times \left( \frac{1}{4} \right) ^n$:}{Write $P_n$ in the form $P_n = \underline{\phantom{4096}} \times \left( \frac{1}{4} \right) ^n$:}

    \vspace{1cm}

  \[
    P_n = \hspace{6cm}
  \]


\end{enumerate}

\newpage

% ════════════════════════════════════════════════════════════════════════════
\sectionbox{\ealpara{Sekcja B}{Section B}}{\ealpara{Różne przedstawienia \ealkeytr{ciągów}}{Different representations of \ealkey{sequences}}}
% ════════════════════════════════════════════════════════════════════════════

\vspace{6pt}
\ealpara{Każdy wiersz poniżej pokazuje ten sam \ealkeytr{ciąg} w trzech różnych formach. Jedna komórka w każdym wierszu jest już wypełniona --- uzupełnij resztę.}{Each row below shows the same \ealkey{sequence} in three different forms. One cell in each row is already filled in --- complete the rest.}

\scaffold{\ealpara{Znasz już niektóre metody znajdowania \ealkeytr{wzoru na $n$-ty wyraz} \ealkeytr{ciągów}, jednak jeśli \ealkeytr{ciąg} wykorzystuje mnożenie, rozważ swoją odpowiedź do {\bf A4} i wcześniejszą pracę o wzroście składanym. Użyj podstawienia, aby sprawdzić swoje odpowiedzi. Uważaj, niektóre odpowiedzi będą używać innej \ealkeytr{zmiennej} niż $n$ we wzorze na ``$n$-ty'' \ealkeytr{wyraz}!}{You already know some methods to find the \ealkey{$n$th term} of \ealkey{sequences}, however if the \ealkey{sequence} uses multiplication then consider you answer to {\bf A4}, and previous work on compounding growth. Use substitution to check your answers. Be careful, some answers will use a different \ealkey{variable} to $n$ in the ``$n$th'' \ealkey{term}!}}
\vspace{12pt}

% ── B1 ──────────────────────────────────────────────────────────────────────
\noindent\textbf{B1.}
\vspace{5pt}

\ealpara{\ealkeytr{Proces iteracyjny} \textbullet\ Tabela (pierwsze 5 \ealkeytr{wyrazów}) \textbullet\ \ealkeytr{Wzór na $n$-ty wyraz}}{\ealkey{Iterative Process} \textbullet\ Table (first 5 \ealkey{terms}) \textbullet\ \ealkey{$n$th Term} Formula}

\noindent
\renewcommand{\arraystretch}{1.3}
\begin{tabularx}{\linewidth}{|>{\centering\arraybackslash}X
                             |>{\centering\arraybackslash}X
                             |>{\centering\arraybackslash}X|}
\hline
\rowcolor{sectionbg}
\textbf{\ealkey{Iterative Process}} & \textbf{Table (first 5 \ealkey{terms})} & \textbf{\ealkey{$n$th Term} Formula} \\
\hline
& \vspace{4pt}$5,\ 8,\ 11,\ 14,\ 17\ldots$\vspace{4pt} & \\[4.5cm]
\hline
\end{tabularx}

\vspace{16pt}

% ── B2 ──────────────────────────────────────────────────────────────────────
\noindent\textbf{B2.}
\vspace{5pt}

\ealpara{\ealkeytr{Proces iteracyjny} \textbullet\ Tabela (pierwsze 5 \ealkeytr{wyrazów}) \textbullet\ \ealkeytr{Wzór na $n$-ty wyraz}}{\ealkey{Iterative Process} \textbullet\ Table (first 5 \ealkey{terms}) \textbullet\ \ealkey{$n$th Term} Formula}

\noindent
\begin{tabularx}{\linewidth}{|>{\centering\arraybackslash}X
                             |>{\centering\arraybackslash}X
                             |>{\centering\arraybackslash}X|}
\hline
\rowcolor{sectionbg}
\textbf{\ealkey{Iterative Process}} & \textbf{Table (first 5 \ealkey{terms})} & \textbf{\ealkey{$n$th Term} Formula} \\
\hline
\vspace{4pt}$F_{j+1} = \dfrac{1}{2}F_j,\quad F_1 = 80$\vspace{4pt} & & \\[4.5cm]
\hline
\end{tabularx}

\vspace{16pt}

\newpage{}

% ── B3 ──────────────────────────────────────────────────────────────────────
\ealpara{\textbf{B3.}\quad \textit{Pewna wielkość maleje o 15\% na każdym kroku, zaczynając od $V_0 = 200$.}}{\textbf{B3.}\quad \textit{A quantity decreases by 15\% each step, starting at $V_0 = 200$.}}

\vspace{5pt}

\ealpara{\ealkeytr{Proces iteracyjny} \textbullet\ Tabela ($n=0$ do $4$, 2 miejsca dziesiętne) \textbullet\ \ealkeytr{Wzór na $n$-ty wyraz}}{\ealkey{Iterative Process} \textbullet\ Table ($n=0$ to $4$, 2 d.p.) \textbullet\ \ealkey{$n$th Term} Formula}

\noindent
\begin{tabularx}{\linewidth}{|>{\centering\arraybackslash}X
                             |>{\centering\arraybackslash}X
                             |>{\centering\arraybackslash}X|}
\hline
\rowcolor{sectionbg}
\textbf{\ealkey{Iterative Process}} & \textbf{Table ($n=0$ to $4$, 2 d.p.)} & \textbf{\ealkey{$n$th Term} Formula} \\
\hline
& & \\[4.5cm]
\hline
\end{tabularx}

\vspace{16pt}

% ── B4 ──────────────────────────────────────────────────────────────────────
\noindent\textbf{B4.}
\vspace{5pt}

\ealpara{W B4 nie istnieje zgrabny wzór --- zobacz pytanie poniżej}{In B4, no neat formula exists --- see question below}

\ealpara{\ealkeytr{Proces iteracyjny} \textbullet\ Tabela (pierwsze 5 \ealkeytr{wyrazów}, 2 miejsca dziesiętne) \textbullet\ \ealkeytr{Wzór na $n$-ty wyraz}}{\ealkey{Iterative Process} \textbullet\ Table (first 5 \ealkey{terms}, 2 d.p.) \textbullet\ \ealkey{$n$th Term} Formula}

\noindent
\begin{tabularx}{\linewidth}{|>{\centering\arraybackslash}X
                             |>{\centering\arraybackslash}X
                             |>{\centering\arraybackslash}X|}
\hline
\rowcolor{sectionbg}
\textbf{\ealkey{Iterative Process}} & \textbf{Table (first 5 \ealkey{terms}, 2 d.p.)} & \textbf{\ealkey{$n$th Term} Formula} \\
\hline
\vspace{4pt}$Q_{n+1} = \dfrac{Q_n}{2} + 6,\quad Q_1 = 100$\vspace{4pt} & & \textit{No neat formula exists --- see question below} \\[4.5cm]
\hline
\end{tabularx}

\vspace{8pt}
\ealpara{Co jest inne w tym pytaniu?}{What is different about this question?}

\answerlines{3}

\newpage

% ════════════════════════════════════════════════════════════════════════════
\sectionbox{\ealpara{Sekcja C}{Section C}}{\ealpara{Zadania kontekstowe --- Kiedy zaczyna się \ealkeytr{ciąg}?}{Context Problems --- When Does the \ealkey{Sequence} Start?}}
% ════════════════════════════════════════════════════════════════════════════

\vspace{4pt}
\ealpara{W tych pytaniach musisz zdecydować, czy \ealkeytr{ciąg} powinien zaczynać się od $n = 0$, czy od $n = 1$, i uzasadnić swój wybór.}{In these questions you must decide whether the \ealkey{sequence} should start at $n = 0$ or $n = 1$, and justify your choice.}

\bigskip

% ── C1 ──────────────────────────────────────────────────────────────────────
\ealpara{\textbf{C1 --- Konto oszczędnościowe}}{\textbf{C1 --- Savings Account}}

\begin{tcolorbox}[colback=white, colframe=lightgrey, sharp corners, left=8pt, top=6pt, bottom=6pt]
\ealpara{Amara otwiera konto oszczędnościowe z \pounds600. Co roku na końcu roku dodawane jest 5\% odsetek. Niech $V_k$ będzie wartością konta \textbf{po $k$ pełnych latach}.}{Amara opens a savings account with \pounds600. Each year, 5\% interest is added at the end of the year. Let $V_k$ be the value of the account \textbf{after $k$ complete years}.}
\end{tcolorbox}

\scaffold{\ealpara{Odsetki są dodawane \emph{na końcu} każdego roku. Zanim dodano jakiekolwiek odsetki (w roku 0), konto zawiera pierwotne \pounds600. Więc ile wynosi $V_0$? Co reprezentuje $V_1$?}{Interest is added \emph{at the end} of each year. Before any interest has been added (at year 0), the account holds the original \pounds600. So what is $V_0$? What does $V_1$ represent?}}

\begin{enumerate}
  \item \ealpara{Podaj wartość $V_0$ i wyjaśnij, dlaczego \ealkeytr{ciąg} zaczyna się od $k = 0$, a nie od $k = 1$.}{State the value of $V_0$ and explain why the \ealkey{sequence} starts at $k = 0$ rather than $k = 1$.}

  \answerlines{2}

  \item \ealpara{Zapisz \ealkeytr{proces iteracyjny} dla $V_{k+1}$ za pomocą $V_k$.}{Write an \ealkey{iterative process} for $V_{k+1}$ in terms of $V_k$.}
  \[
    V_{k+1} = \hspace{6cm}, \qquad V_0 = \hspace{2cm}
  \]

  \item \ealpara{Uzupełnij tabelę dla $k = 0$ do $4$. Podaj wartości z dokładnością do jednego pensa.}{Complete the table for $k = 0$ to $4$. Give values to the nearest penny.}

  \vspace{6pt}
  \begin{center}
  \renewcommand{\arraystretch}{1.8}
  \begin{tabular}{|>{\bfseries\centering\arraybackslash}m{1.2cm}
                  |>{\centering\arraybackslash}m{2.0cm}
                  |>{\centering\arraybackslash}m{2.0cm}
                  |>{\centering\arraybackslash}m{2.0cm}
                  |>{\centering\arraybackslash}m{2.0cm}
                  |>{\centering\arraybackslash}m{2.0cm}|}
  \hline
  \rowcolor{sectionbg}
  $k$        & 0        & 1 & 2 & 3 & 4 \\
  \hline
  $V_k$ (\pounds) & \pounds600.00 &   &   &   &   \\
  \hline
  \end{tabular}
  \end{center}

  \vspace{8pt}
  \item \ealpara{Zapisz wzór na $V_k$ bezpośrednio. Użyj go, aby obliczyć $V_{10}$ (z dokładnością do jednego pensa).}{Write a formula for $V_k$ directly. Use it to calculate $V_{10}$ (to the nearest penny).}

  \vspace{4pt}
  $V_k = \hspace{7cm}$

  \vspace{0.5cm}
  $V_{10} = \hspace{7cm}$
  \vspace{4pt}

\end{enumerate}

\bigskip\bigskip

\newpage{}

% ── C2 ──────────────────────────────────────────────────────────────────────
\ealpara{\textbf{C2 --- Piłka odbijająca się}}{\textbf{C2 --- Bouncing Ball}}

\begin{tcolorbox}[colback=white, colframe=lightgrey, sharp corners, left=8pt, top=6pt, bottom=6pt]
\ealpara{Piłkę upuszczono z wysokości 4\,m. Każde odbicie osiąga 75\% poprzedniej wysokości.\\[4pt]
\textbf{Jess} definiuje $H_n$ jako wysokość osiągniętą po $n$ odbiciach, pomijając początkowy upadek.\\
\textbf{Kwame} definiuje $H_n$ jako wysokość po $n$ odbiciach, gdzie \emph{początkowa wysokość upadku liczy się jako wysokość po 0 odbiciach}.}{A ball is dropped from a height of 4\,m. Each bounce reaches 75\% of the previous height.\\[4pt]
\textbf{Jess} defines $H_n$ as the height reached after $n$ bounces ignoring the initial drop.\\
\textbf{Kwame} defines $H_n$ as the height after $n$ bounces, where the \emph{original drop height counts as the height after 0 bounces}.}
\end{tcolorbox}


\begin{enumerate}
  \item \ealpara{Zapisz \ealkeytr{proces iteracyjny} Jess i podaj jej wartość początkową.}{Write Jess's \ealkey{iterative process} and state her starting value.}
  \[
    H_{n+1} = \hspace{5cm}, \qquad H_1 = \hspace{2cm}
  \]

  \item \ealpara{Zapisz \ealkeytr{proces iteracyjny} Kwame'a i podaj jego wartość początkową.}{Write Kwame's \ealkey{iterative process} and state his starting value.}
  \[
    H_{n+1} = \hspace{5cm}, \qquad H_0 = \hspace{2cm}
  \]

  \item \ealpara{Co jest takie samo w \ealkeytr{procesach iteracyjnych} Jess i Kwame'a?}{What is the same about Jess's and Kwame's \ealkey{iterative processes}?}

  \answerlines{2}

  \item \ealpara{Co jest inne w \ealkeytr{procesach iteracyjnych} Jess i Kwame'a?}{What is different about Jess's and Kwame's \ealkey{iterative processes}?}

  \answerlines{2}


  \item \ealpara{Dla \textbf{każdej} definicji podaj, jaką wysokość reprezentuje $H_3$, i oblicz jej wartość.}{For \textbf{each} definition, state what height $H_3$ represents, and calculate its value.}

  \vspace{6pt}
  \ealpara{\textit{Jess:}\quad $H_3 = $ \underline{\hspace{2cm}} m. \quad To oznacza: \underline{\hspace{7cm}}}{\textit{Jess:}\quad $H_3 = $ \underline{\hspace{2cm}} m. \quad This represents: \underline{\hspace{7cm}}}

  \vspace{10pt}
  \ealpara{\textit{Kwame:}\quad $H_3 = $ \underline{\hspace{2cm}} m. \quad To oznacza: \underline{\hspace{6.15cm}}}{\textit{Kwame:}\quad $H_3 = $ \underline{\hspace{2cm}} m. \quad This represents: \underline{\hspace{6.15cm}}}

  \vspace{10pt}
  \item \ealpara{Czyja definicja bardziej ci odpowiada?}{Whose definition do you prefer?}

  \answerlines{3}

\end{enumerate}

\newpage

% ════════════════════════════════════════════════════════════════════════════
% ════════════════════════════════════════════════════════════════════════════
\sectionbox{\ealpara{Sekcja D}{Section D}}{\ealpara{Zadania rozszerzające i wyzwania}{Extension and Challenge Problems}}
% ════════════════════════════════════════════════════════════════════════════


\bigskip

% ── D1 ──────────────────────────────────────────────────────────────────────
\ealpara{\textbf{D1 --- Praca wstecz}}{\textbf{D1 --- Working Backwards}}

\begin{tcolorbox}[colback=white, colframe=lightgrey, sharp corners, left=8pt, top=6pt, bottom=6pt]
\ealpara{5. \ealkeytr{wyraz} pewnego iteracyjnego \ealkeytr{ciągu} to $U_5 = 48$.\\
\ealkeytr{Proces iteracyjny} to \quad $U_{n+1} = 2U_n$.}{The 5th \ealkey{term} of an iterative \ealkey{sequence} is $U_5 = 48$.\\
The \ealkey{iterative process} is \quad $U_{n+1} = 2U_n$.}
\end{tcolorbox}

\begin{enumerate}
  \item \ealpara{Znajdź $U_1$. Pokaż wszystkie kroki jasno.}{Find $U_1$. Show all steps clearly.}

  \answerlines{4}

  \item \ealpara{Zapisz \ealkeytr{wzór na $n$-ty wyraz} dla $U_n$. Użyj go, aby sprawdzić, że $U_5 = 48$.}{Write an \ealkey{$n$th term} formula for $U_n$. Use it to verify that $U_5 = 48$.}

  \answerlines{3}

\end{enumerate}

\bigskip\bigskip

\newpage{}

% ── D2 ──────────────────────────────────────────────────────────────────────
\ealpara{\textbf{D2 --- \ealkeytr{Zbieżność}}}{\textbf{D2 --- \ealkey{Convergence}}}

\begin{tcolorbox}[colback=white, colframe=lightgrey, sharp corners, left=8pt, top=6pt, bottom=6pt]
\ealpara{Rozważ \ealkeytr{proces iteracyjny} \quad $F_{j+1} = \dfrac{F_j}{2} + 6, \qquad F_1 = 100$}{Consider the \ealkey{iterative process} \quad $F_{j+1} = \dfrac{F_j}{2} + 6, \qquad F_1 = 100$}
\end{tcolorbox}

\begin{enumerate}
  \item \ealpara{Uzupełnij tabelę z dokładnością do 2 miejsc dziesiętnych.}{Complete the table to 2 decimal places.}

  \vspace{6pt}
  \begin{center}
  \renewcommand{\arraystretch}{1.9}
  \begin{tabular}{|>{\bfseries\centering\arraybackslash}m{1.2cm}
                  |>{\centering\arraybackslash}m{1.5cm}
                  |>{\centering\arraybackslash}m{1.5cm}
                  |>{\centering\arraybackslash}m{1.5cm}
                  |>{\centering\arraybackslash}m{1.5cm}
                  |>{\centering\arraybackslash}m{1.5cm}
                  |>{\centering\arraybackslash}m{1.5cm}
                  |>{\centering\arraybackslash}m{1.5cm}
                  |>{\centering\arraybackslash}m{1.5cm}|}
  \hline
  \rowcolor{sectionbg}
  $j$   & 1   & 2 & 3 & 4 & 5 & 6 & 7 & 8 \\
  \hline
  $F_j$ & 100 &   &   &   &   &   &   &   \\
  \hline
  \end{tabular}
  \end{center}

  \vspace{8pt}
  \item \ealpara{Do jakiej wartości $F_j$ wydaje się zbliżać?}{What value does $F_j$ appear to be approaching?}

  \answerlines{1}

  \item \ealpara{Gdyby \ealkeytr{ciąg} ostatecznie ustabilizował się na stałej wartości $L$ (zwanej \ealkeytr{punktem stałym}), spełniałby}{If the \ealkey{sequence} eventually settled at a fixed value $L$ (called a \ealkey{fixed point}), it would satisfy}
  \[
    L = \frac{L}{2} + 6
  \]
  \ealpara{Rozwiąż to równanie, aby znaleźć $L$.}{Solve this equation to find $L$.}

  \answerlines{3}

  \item \ealpara{Co zauważasz w swoich odpowiedziach do (b) i (c)? Czy uważasz, że $F_j$ kiedykolwiek będzie miał \ealkeytr{wyraz}, który ma tę samą wartość co $L$?}{What do you notice about your answers to (b) and (c)? Do you think $F_j$ will ever have a \ealkey{term} that has the same value as $L$?}

  \answerlines{3}

\end{enumerate}

\bigskip\bigskip

\end{document}